\section{应用案例与演化时序}

\subsection{Benchmark 背景与时间线}

\begin{knowledgebox}{Q-Learning 演化时间线}
\begin{itemize}
    \item 1992 年：Watkins 的 Q-Learning 理论收敛。
    \item 2013 年：Mnih 等在 Atari 上用 DQN 实现 end-to-end 控制。
    \item 2015-2017 年：Double DQN、Dueling、Prioritized Replay、Rainbow 等陆续完善。
    \item 2016 年后：DQN 拓展到 MuJoCo、Roboschool 等连续控制任务，多 agent RL 与对抗训练开始流行。
\end{itemize}
\end{knowledgebox}

提出 benchmark 时总会面对一个抉择：是用简单环境获取重复性、还是用复杂环境逼近真实应用。Stanford 的课程复习强调了前者的控制性与后者的代表性之间的 trade-off。

\subsection{案例洞察}

\begin{importantbox}{Atari 案例复盘}
在 Atari Breakout 任务上，Agent 通过 frame stacking 提取动量信息，而 frame-skip 提高了训练效率。改变 reward shaping（例如加大 brick 击中奖励）后，agent 的策略发生了明显偏移，说明 reward scales 的调整容易触发新策略。
\end{importantbox}

\begin{warningbox}{现实应用中的观测偏差}
在机器人任务中，真实 sensor 噪声和模拟 sensor 差异可能导致 sim-to-real 的失败。即便在仿真中训练到高分，也必须做 domain randomization 并在部署前进行实际数据校对。
\end{warningbox}

\subsection{本章小结}

本章结合 benchmark 发展与经典案例，提醒我们即便掌握了算法细节，也要当心仿真与真实世界之间的信息缺口。

